\section{Conclusion}\label{sec:conclusion}

This study examined a series of probabilistic models to investigate correlations between factors of wildfires and their surrounding communities and the order of an evacuation. The change from positive to negative coefficient on the FN indicator variable after introducing further covariates suggests that FN reserves are correlated with a lower likelihood of an evacuation order. The stability of this conclusion after giving the model access to response severity and controlling for `remoteness' using a grid of regions adds weight to the claim. This could suggest bias in the decision-making process fuelled by priorities embedded in settler systems of power, but the results observed here should not be taken in isolation as suggesting a causal relationship. Other statistical methods and qualitative analysis of such systems should be used to investigate causality.

\textbf{Limitations:} Among the main limitations is the data. The join between fires and evacuations is imperfect, and the removal of unmatched evacuations may bias results. Further, the data is observational and thus subject to confounding and bias. Ideas around `remoteness' are hard to capture, even by local regions, and may be a meaningful confounder. The model is also limited in its ability to capture the complex decision-making process around evacuation, which may involve many unobserved factors. Further, not all fires are recorded in the data, and not all evacuations are recorded. A `hurdle' model that accounts for the fact that some fires are not reported may be a useful extension along with identifying higher-order effects or interaction terms.